\documentclass[10pt,a4paper]{amsart}
\usepackage[margin=19mm]{geometry}
\usepackage{amsmath,amssymb,mathtools}
\usepackage[scr=rsfs,cal=euler]{mathalfa}
\usepackage{xfrac}
\usepackage[unicode,hidelinks,bookmarks=false]{hyperref}
\newcommand{\ie}{i.e.\ }
\newcommand{\cf}{cf.\ }
\newcommand{\resp}{respectively\ }
\newcommand{\Subsection}{\subsection}
\providecommand{\nobreakdash}{\nobreak\hbox{-}}
\setlength{\emergencystretch}{2em}
\multlinegap0pt
\usepackage{bm}
\usepackage{bbm}
\usepackage{dsfont}
\usepackage{xspace}
\usepackage{cancel}
\usepackage{mathtools}
\usepackage{amsthm}
\makeatletter
\@ifundefined{theorem}{\newtheorem{theorem}{Theorem}}{}
\@ifundefined{lemma}{\newtheorem{lemma}[theorem]{Lemma}}{}
\@ifundefined{proposition}{\newtheorem{proposition}[theorem]{Proposition}}{}
\makeatother
\makeatletter
\newcommand{\bR}{\mathbb{R}}
\newcommand{\bZ}{\mathbb{Z}}
\expandafter\ifx\csname blue\endcsname\relax
\newenvironment{blue}{\relax\color{blue}}{\hspace*{.5ex}\relax}
\fi
\newcommand{\cS}{\mathcal{S}}
\expandafter\ifx\csname red\endcsname\relax
\newenvironment{red}{\relax\color{red}}{\hspace*{.5ex}\relax}
\fi
\newcommand{\wtilde}{\widetilde}
\makeatother
\title[Positive quasimodular forms and the sign uncertainty princip]{Positive quasimodular forms and the sign uncertainty principle}
\author{Seewoo Lee}
\date{Source: arXiv:2608.15415v2 (CC BY 4.0)}
\begin{document}
\maketitle

\noindent\emph{Source and licence.} Positive quasimodular forms and the sign uncertainty principle. Version arXiv:2608.15415v2 (CC BY 4.0), distributed under the Creative Commons Attribution 4.0 International licence, as recorded in the arXiv licence metadata for this exact version and shown on the abstract page https://arxiv.org/abs/2608.15415v2. Full rights notice: licenses/arxiv-2608.15415-CC-BY-4.0.txt.\par\smallskip

\noindent\textit{Editorial scope.} Counted unit: the Proposition whose source label is \texttt{prop:Ftilde\_pos} in the article (source0.tex lines 1488--1494, source label \texttt{prop:Ftilde\_pos}), delivered together with the article's own complete original proof of it (source0.tex lines 1495--1543, one \texttt{proof} environment of 49 lines). No mathematics is written, repaired or re-derived: statement and proof are the article's own text line for line. Required context retained uncounted: lem:KZ2\_coeff (lines 895--908); eqn:Ftilde\_rec (lines 1466--1469). Every definition, display and label of those lines is retained unchanged. Omitted from this package and not redistributed: 1--894, 909--1465, 1470--1487, 1544--2420. Nothing inside a retained excerpt is omitted or altered: every retained body is delivered exactly as the article's lines are written, so no omission receipt of any kind accompanies this package. Citations: the retained excerpts carry no citation command at all, so no bibliography entry of the article is transcribed and no reference is redistributed. Reference adaptation: none was needed and none was made; every reference inside the delivered unit resolves inside it, so no unresolved reference or citation is produced. Printed slips kept literal: no printed word, symbol, hypothesis or number of the counted statement or of its proof was altered. Layout and class: the article's own class is \texttt{amsproc} with packages including fullpage, xcolor, inputenc, amsmath, amsthm, amssymb, mathtools, dsfont, units, graphicx, wrapfig, subfig, float, listings; the wrapper is the lane's pinned \texttt{amsart} editorial wrapper at 10pt on A4, which restates the theorem-like environments the retained text needs and adds the article's own macro definitions that the retained text needs (\texttt{\textbackslash bR}, \texttt{\textbackslash bZ}, \texttt{\textbackslash cS}, \texttt{\textbackslash wtilde}), with extra wrapper packages (bm, bbm, dsfont, xspace, cancel, mathtools). The delivered numbering is the unit's own. Assets: no retained span contains a figure environment, a graphics-inclusion command, a PDF-page inclusion, a table or a TikZ picture; no image file is copied into this package, the package declares no asset dependency and no image file is redistributed with it. Licence: version arXiv:2608.15415v2 of this article is distributed under the Creative Commons Attribution 4.0 International licence, as recorded in the arXiv licence metadata for this exact version and shown on the abstract page https://arxiv.org/abs/2608.15415v2. A full rights notice is delivered at licenses/arxiv-2608.15415-CC-BY-4.0.txt. Characters: every retained excerpt is pure ASCII, so the delivered engine, XeLaTeX with its default Latin Modern text font, reports no missing character.
\begin{lemma}
    \label{lem:KZ2_coeff}
    Let $G = \sum_{n \ge 0} a_n q^n$ be a quasimodular form, and let $n \ge 0$.
    For $k \in \bZ$ and $\alpha \in \bR$, the $n$-th Fourier coefficient of $L_{2,k}^{\alpha} G$ is
    \begin{equation}
        \label{eqn:KZ2_coeff}
        [q^n] (L_{2,k}^{\alpha} G) = \kappa_{2,k}^{\alpha}(n) a_n + \sum_{j=0}^{n-1} K_{2,k}^{\alpha}(n,j) a_j,
    \end{equation}
    where
    \begin{align}
        \kappa_{2,k}^{\alpha}(n) &= n^2 - \frac{k+1}{6}n + \alpha, \label{eqn:kappa2} \\
        K_{2,k}^{\alpha}(n,j) &= 2(k+1)(2j-k(n-j)) \sigma_1(n-j) + 240 \alpha \sigma_3(n-j). \label{eqn:K2}
    \end{align}
\end{lemma}

    \begin{equation}
        \label{eqn:Ftilde_rec}
        \widetilde{F}_{w+2} = \frac{3(w-4)w}{16(w-10)(w-5)(w-3)(w+2)} \left(\mathcal{S}_w \widetilde{F}_{w-2} - \frac{1}{3} \partial_{w-1} F_{w}\right)
    \end{equation}

\begin{proposition}
    \label{prop:Ftilde_pos}
    For $w \equiv 0 \pmod{4}$ and $w \ge 12$, we have 
    \begin{equation}
        \tilde{a}_{j,+}^{(w-2)} > 0 \quad \text{for all} \quad 1 \le j \le \frac{w}{4} - 3, \quad \tilde{a}_{\frac{w}{4}-2,+}^{(w-2)} \ge \frac{1}{360}.
    \end{equation}
\end{proposition}

\begin{proof}
    Let $\alpha_w = - \frac{(w-10)(w-4)}{48}$.
    From $\cS_w = -L_{2,w-2}^{\alpha_w}$, Lemma \ref{lem:KZ2_coeff} implies that the $n$-th Fourier coefficient of $\cS_w \wtilde{F}_{w-2}$ is
    \[
        [q^n](\cS_w \wtilde{F}_{w-2}) = -\kappa_{2,w-2}^{\alpha_w}(n) \tilde{a}_{n,+}^{(w-2)} - \sum_{j=1}^{n-1} K_{2,w-2}^{\alpha_w}(n, j) \tilde{a}_{j,+}^{(w-2)}
    \]
    where
    \begin{align}
        \kappa_{2,w-2}^{\alpha_w}(n) &= \frac{(12n + w - 10)(4n - w + 4)}{48} \label{eqn:kappa_F}, \\
        K_{2,w-2}^{\alpha_w}(n,j) &= -2(w-1)\bigl(w(n-j)-2n\bigr)\sigma_1(n-j) - 5(w-10)(w-4)\sigma_3(n-j). \label{eqn:K_F}
    \end{align}
    We now use induction on $w$.
    For $w = 12$, direct computation shows that
    \[
    \wtilde{F}_{10} = \frac{1}{360} E_4 X_{6,1}
    \]
    which is completely positive, i.e. $\tilde{a}_{j,+}^{(10)} > 0$ for all $j \ge 1$.
    Also, $\tilde{a}_{1,+}^{(10)} = \frac{1}{360}$.
    Assume that $\tilde{a}_{j,+}^{(w-2)} > 0$ for all $1 \le j \le \frac{w}{4} - 3$ and $\tilde{a}_{\frac{w}{4}-2,+}^{(w-2)} \ge \frac{1}{360}$.
    Then, by \eqref{eqn:Ftilde_rec}, \eqref{eqn:kappa_F}, and \eqref{eqn:K_F}, the $n$-th Fourier coefficient of $\wtilde{F}_{w+2}$ is
    \begin{equation}
        \label{eqn:Ftilde_induction_coeff}
        c_w\left(-\kappa_{2,w-2}^{\alpha_w}(n) \tilde{a}_{n,+}^{(w-2)} - \sum_{j=1}^{n-1} K_{2,w-2}^{\alpha_w}(n,j)\tilde{a}_{j,+}^{(w-2)} - \frac{1}{3}n b_n^{(w)} + \frac{w-1}{36}\left(b_n^{(w)} - 24 \sum_{k=1}^{n-1} \sigma_1(n-k) b_k^{(w)}\right)\right)
    \end{equation}
    where $c_w = \frac{3(w-4)w}{16(w-10)(w-5)(w-3)(w+2)}>0$.
    Since $F_w = q^{\frac{w}{4} - 1} + O(q^{\frac{w}{4}})$, we have $b_n^{(w)} = 0$ for all $1 \leq n \leq \frac{w}{4} - 2$.
    In this range, $4n-w+4<0$, so \eqref{eqn:kappa_F} gives $\kappa_{2,w-2}^{\alpha_w}(n)<0$. Moreover, for $1\leq j\leq n-1$ and $n\leq\frac{w}{4}-1$,
    \[
        w(n-j)-2n \geq w-2n>0,
    \]
    and hence \eqref{eqn:K_F} gives $K_{2,w-2}^{\alpha_w}(n,j)<0$.
    For such $n$, \eqref{eqn:Ftilde_induction_coeff} becomes
    \[
        c_w\left(- \kappa_{2,w-2}^{\alpha_w}(n) \tilde{a}_{n,+}^{(w-2)} - \sum_{j=1}^{n-1} K_{2,w-2}^{\alpha_w}(n,j)\tilde{a}_{j,+}^{(w-2)}\right)
    \]
    which is positive by the induction hypothesis.
    For $n = \frac{w}{4} - 1$, normalization gives $b_n^{(w)}=1$ and $b_k^{(w)}=0$ for $k<n$, while $\kappa_{2,w-2}^{\alpha_w}(n) = 0$ and
    \[
    K_{2,w-2}^{\alpha_w}\left(\frac{w}{4}-1, \frac{w}{4}-2\right) = -2(w-1)\left(w - 2\left(\frac{w}{4}-1\right)\right) - 5(w-10)(w-4) = -(6w^2 - 67w + 196),
    \]
    hence \eqref{eqn:Ftilde_induction_coeff} can be bounded from below as
    \begin{align}
        c_w\left(-\sum_{j=1}^{n-1} K_{2,w-2}^{\alpha_w}(n,j)\tilde{a}_{j,+}^{(w-2)} - \frac{2w-11}{36}\right) &\ge c_w \left(-K_{2,w-2}^{\alpha_w}\left(\frac{w}{4}-1, \frac{w}{4}-2\right) \tilde{a}_{\frac{w}{4}-2,+}^{(w-2)} - \frac{2w-11}{36}\right) \nonumber \\
        &\ge c_w \left(\frac{6w^2 - 67w + 196}{360} - \frac{2w - 11}{36}\right) \nonumber \\
        &= \frac{(w-6)(w-4)w(2w-17)}{640(w-10)(w-5)(w-3)(w+2)} \label{eqn:boundary_coeff}
    \end{align}
    and \eqref{eqn:boundary_coeff} is greater than or equal to $\frac{1}{360}$ if and only if $2w^4 - 77w^3 + 1018 w^2 - 4312 w + 4800 \ge 0$, which is true for all $w \ge 12$.
    This proves the claim for $n = \frac{w}{4} - 1$ and completes the induction.
\end{proof}

\end{document}
